\section{Conclusion}
\label{sec:conclusion}

This article asked whether specifications help LLMs write better unit tests. Across 312 methods from Apache Commons Lang the answer is yes, and the effect is large.

\begin{enumerate}
    \item \emph{Specifications substantially improve test quality.} LLM-generated tests guided by specifications produce 272\% more tests and achieve a mutation score 40.7~pp higher than signature-only baselines ($p < 0.001$, Cohen's $d = 2.41$, 95\% bootstrap CI $[2.08, 2.79]$).

    \item \emph{The effect varies systematically across method categories.} Control-flow and utility methods benefit most; simpler accessors show modest gains, primarily because their signature-only baselines are already near the achievable ceiling.

    \item \emph{Specifications close most, but not all, of the gap to source code.} Full source access (P4) attains the highest absolute mutation score (94.8\%); the specification phase reaches 68.2\%, well above the signature-only baseline of 27.5\%, while remaining materially more compact than full source.
\end{enumerate}

The specifications used in this study were produced automatically by JML-Inferrer, a heuristic static analysis system that emits non-trivial specifications for 99.0\% of the corpus and discharges 92.9\% of inferred clauses through OpenJML's Extended Static Checker. The system is the means by which the empirical question above could be answered at scale, and addresses the principal barrier to specification adoption---manual effort~\cite{zimmermann2019lightweight, hall1990seven}---while remaining complementary to documentation-based approaches~\cite{jdoctor2018}.

The result above is the headline use of inferred specifications, but it is not the only use. The same artefacts function as oracle-bearing input for any test-generation tool that consumes JML; as the callee contracts that compositional verification needs in order to check calling code without re-analysing the entire transitive dependency graph; and as a machine-readable contract that AI coding agents can consume when reasoning about a library, particularly one that postdates the model's training cut-off. The agentic-programming case is itself the source of the development methodology used to build JML-Inferrer: candidate JML produced by an LLM (or by hand) is verified with OpenJML, validated examples become probe tests, and an agentic coding assistant extends the inferrer to make those probes pass. The same workflow is available to users of the tool who need it to recognise a new pattern in their own code.

The findings reported here motivate a planned follow-up study targeting service-boundary code. Methods that call REST endpoints or RPC stubs are precisely the setting in which the oracle problem is most acute: the implementation typically delegates to an external service, signatures convey nothing about status codes or retry semantics, and signature-only LLM test generation collapses into mock-the-call assertions. Whether inferred specifications retain or amplify the gain reported here in that setting is the central question of the planned follow-up. Further directions include strengthening loop invariant inference through machine-learning~\cite{wei2011inferring} or hybrid static--dynamic techniques~\cite{ernst2007daikon}, explicit treatment of object-oriented constructs (inheritance, polymorphism, class invariants), extension to languages beyond Java, interactive refinement under developer guidance~\cite{cousot2011contract}, and recovery of abstract-contract over-approximations in cases where the implementation makes a choice that the API leaves open (sort stability for duplicates being the canonical case)---a class of specification not derivable from a single implementation in isolation. The replication package, including all source code, data, generated test suites, and analysis scripts, accompanies this article.
